%\pagebreak
\section{\kcode{flatten} Construct}
\label{sec:flatten}
\index{constructs!flatten@\kcode{flatten}}
\index{flatten construct@\kcode{flatten} construct}

The \kcode{flatten} construct combines a (canonical) loop nest into a single flattened loop.
The \kcode{depth} clause specifies the number of loops to be flattened. 
When the \kcode{depth} clause is not present or no argument is specified,
the default of 2 is used.  Its basic operation is similar to the \kcode{collapse} 
clause used with work-sharing/work-distribution and parallelism generating directives.
However, the \kcode{flatten} construct and its \kcode{apply} clause allow 
it to be applied in a more generic way to loop-transforming constructs. 
Also, flattening can be applied to non-rectangular loops.

In the following example the \kcode{flatten} construct is applied to the outer two loops,
in the \ucode{flatten} routine.  The innermost \ucode{k} loop is not affected.

The \ucode{flatten_equivalent} routine is an illustration of how a user might 
write simple code that flattens the loops.
The specifics of the actual integer type used to represent the loop iteration variable 
for the directive flattened loop is implementation defined.
For instance, in C/C++ if \ucode{Ni} and \ucode{Nj} are \bcode{int} type and the loop 
iteration variable ranges are positive integers, the compiler may choose 
\bcode{unsigned int} to represent the loop iteration
\ucode{ij} and trip count \ucode{tc_end} variables.

\cexample[6.1]{flatten}{1}
\ffreeexample[6.1]{flatten}{1} 

The following example is to unroll flattened loops for instruction cache efficiency.
Since the \kcode{unroll} construct only applies to a single loop, 
the \kcode{flatten} construct is used to form a single loop.
The \ucode{flatten_equivalent} routine shows how a user might code the single unrolled loop.

\cexample[6.1]{flatten}{2}
\ffreeexample[6.1]{flatten}{2} 

%%%%%%%%%%%%%%%%%%%%%%%%%%%%%%%%%%%%%%%%%%%%%%%%%%%%%%%%%%%%%%%%%%%%%%%%%%%%%%%%%%%%%
%In a OpenMP loop nest, a non-rectangular loop has a loop bound that references the iteration
%variable of an enclosing loop.

The following example is to flatten a non-rectangular loop nest in which 
the number of iterations in the inner loops depends on the value of 
the outer loop iteration variable.

The \ucode{flatten_equivalent_empty_iterations} routine illustrates
how empty iterations are conditionally skipped when iterating
over a simple direct product iteration space (\ucode{n*n}).
Conditional checks and other methods 
%(such as \texttt{recount}, \texttt{precompute} and \texttt{isqrt})
might be used exclusively or in combination in compiler 
optimization when flattening loops.


\cexample[6.1]{flatten}{3}
\ffreeexample[6.1]{flatten}{3} 

In this example the flattened loops include a non-rectangular loop, as 
above, and the resulting loop is made into a worksharing loop.

Note, the flattened loop formed with a non-rectangular loop may use a 
\kcode{schedule} clause with any schedule type.  This is unlike the 
use of a \kcode{collapse} clause on a worksharing-loop construct that cannot 
use a \kcode{schedule} clause when one of the affected loops is a 
non-rectangular loop.

\cexample[6.1]{flatten}{4}
\ffreeexample[6.1]{flatten}{4} 
